% === Begin Label Data ===
% ID: 166
% 难度星级: 
% 标签: 
% 备注: 
% 组卷引用次数: 0
% === End  Label Data ===

\begin{problem}{2010}{G}{江西卷（理）}{22}{数列}
证明以下命题：

(1) 对任意正整数 $a$，都存在正整数 $b,c\ (b<c)$ 使得 $a^2, b^2, c^2$ 为等差数列.

(2) 存在无穷多互不相似的三角形 $\triangle_n$，其边长 $a_n, b_n, c_n$ 为正整数且 $a_n^2, b_n^2, c_n^2$ 成等差数列.
\end{problem}

\begin{solutions}
第一问

根据题意，正整数 $a,b,c$ 且 $b<c$，满足 $a^2, b^2, c^2$ 为等差数列，即：
$$
b^2-a^2=c^2-b^2 \Rightarrow 2b^2=a^2+c^2
$$
观察上式，我们可以得到 2 个推论：
1. $a<b<c$;
2. $a$ 和 $c$ 要么都是偶数，要么都是奇数。

因为 $a$ 为任意正整数，那么就从最简单的开始，不妨设 $a=1$，则有：
1. 令 $c=1$，得 $b=1$，与 $a<b<c$ 矛盾;
2. 令 $c=3$，得 $b=\sqrt{5}$，与 $b$ 为正整数矛盾;
3. 令 $c=5$，得 $b=\sqrt{13}$，与 $b$ 为正整数矛盾;
4. 令 $c=7$，得 $b=5$，满足条件。

简单猜测实验得到一组要求的值：$a=1, b=5, c=7$。

再设 $a=2$，同理有：
1. 令 $c=4$，得 $b=\sqrt{10}$，与 $b$ 为正整数矛盾;
2. 令 $c=6$，得 $b=2\sqrt{5}$，与 $b$ 为正整数矛盾;
3. 令 $c=8$，得 $b=\sqrt{34}$，与 $b$ 为正整数矛盾;
4. 令 $c=10$，得 $b=2\sqrt{13}$，与 $b$ 为正整数矛盾;
5. 令 $c=12$，得 $b=\sqrt{74}$，与 $b$ 为正整数矛盾;
6. 令 $c=14$，得 $b=10$，满足条件。

通过实验又得到一组要求的值：$a=2, b=10, c=14$。

综合 $a$ 为奇数和偶数的情况，猜想：
对于 $\forall a=n\in \mathbf{N}^+$，令 $b=5a, c=7a$，易得：
$$
2\cdot (5n)^2=n^2+(7n)^2
$$
故命题 (1) 得证。

第二问

相比命题 (1)，我们需要找到无穷多组数 $a,b,c$，必须能构成三角形，且满足 $2b^2=a^2+c^2$。

关于三角形的 3 条边的关系，且满足 $2b^2=a^2+c^2$，很容易想到直角三角形的 3 条边满足：$z^2=x^2+y^2$，这里 $z$ 为直角三角形斜边，$x,y$ 为 2 条直角边。但是我们要找的三角形是 $a<b<c$，如何将两者联系起来呢？

根据直角三角形勾股定理 (Pythagorean Theorem)：$z^2=x^2+y^2$，适当变形：
$$
2z^2=2x^2+2y^2=(y-x)^2+(y+x)^2
$$
上式和题设条件很相似了，令 $a=y-x, b=z, c=x+y$，且 $0<x<y<z$，但 $x,y,z$ 是 3 个未知数，需要 3 个方程才能找到解，而目前只知道 $z^2=x^2+y^2$ 这个关系。

仍然先从最简单的开始，常见的勾股数 (Pythagorean triple) 有：
1. 令 $(x,y,z)$ 为 $(3,4,5)$，则 $(a,b,c)$ 为 $(1,5,7)$，但无法构成三角形;
2. 令 $(x,y,z)$ 为 $(5,12,13)$，则 $(a,b,c)$ 为 $(7,13,17)$，可以构成三角形;
3. 令 $(x,y,z)$ 为 $(7,24,25)$，则 $(a,b,c)$ 为 $(17,25,31)$，可以构成三角形;
4. 令 $(x,y,z)$ 为 $(8,15,17)$，则 $(a,b,c)$ 为 $(9,17,23)$，可以构成三角形;

根据上述分析，出现 2 种情况，$x$ 为奇数和偶数。

$x$ 为奇数

先猜想 $x$ 为奇数时满足条件的表达式为：$(x,y,y+1)$，则有：
$$
x^2+y^2=(y+1)^2 \Rightarrow y=\frac{x^2-1}{2}
$$
设 $x=2n+1(n>2, n\in \mathbf{N}^+)$，则有 $y=2n^2+2n$，故满足题设条件的 $(x,y,z)$ 模式为 $(2n+1, 2n^2+2n, 2n^2+2n+1)$。

易得 $(a,b,c)$ 为
$$
\begin{cases}
a=(2n^2+2n)-(2n+1)=2n^2-1 \\
b=2n^2+2n+1 \\
c=(2n^2+2n)+(2n+1)=2n^2+4n+1
\end{cases}
$$
下面来验证 $(a_n, b_n, c_n)$ 是否可以构成三角形：

当 $\forall n>2$ 时，
$a_n+b_n-c_n=(2n^2-1)+(2n^2+2n+1)-(2n^2+4n+1)=2n(n-1)>0$，故可以构成三角形三边。

下面来证明这些三角形都互不相似，设 $m,n\in \mathbf{N}^+, m\neq n$ 且 $m,n>2$，则有 $a_m=2m^2-1, b_m=2m^2+2m+1, c_m=2m^2+4m+1$; $a_n=2n^2-1, b_n=2n^2+2n+1, c_n=2n^2+4n+1$。

若 $\triangle_m$ 和 $\triangle_n$ 相似，则有：
$$
\frac{a_m}{a_n}=\frac{b_m}{b_n}=\frac{c_m}{c_n}
$$
即：
$$
\frac{2m^2-1}{2n^2-1}=\frac{2m^2+2m+1}{2n^2+2n+1}=\frac{2m^2+4m+1}{2n^2+4n+1}
$$
简单化简则有：
$$
\frac{2n^2+2n+1}{2n}=\frac{2m^2+2m+1}{2m} \Rightarrow 2n+\frac{1}{n}+2=2m+\frac{1}{m}+2
$$
即 $2mn=1$ 而这与假设条件 $m,n>2$ 矛盾，至此我们证明了命题 (2)。

$x$ 为偶数

对于这道题来说，可以不用分析 $x$ 为偶数情况，但这里我们讨论下偶数的情况。满足条件的表达式为：$(x,y,y+2)$，设 $x=2n(n\ge 4, n\in \mathbf{N}^+)$，则满足题设条件的 $(x,y,z)$ 模式为 $(2n, n^2-1, n^2+1)$。

易得 $(a,b,c)$ 为 $(n^2-2n-1, n^2+1, n^2+2n-1)$，下面来验证下 $(a_n, b_n, c_n)$ 是否可以构成三角形：

当 $\forall n \ge 4$ 时，
$a_n+b_n-c_n=(n^2-2n-1)+(n^2+1)-(n^2+2n-1)=n^2-4n+1>0$，故可以构成三角形三边。

同理可设 $m,n\in \mathbf{N}^+, m\neq n$ 且 $m,n \ge 4$，则有：
$$
\frac{m^2-2m-1}{n^2-2n-1}=\frac{m^2+1}{n^2+1}=\frac{m^2+2m-1}{n^2+2n-1}
$$
化简则有：
$$
\frac{m^2-2m-1}{4m}=\frac{n^2-2n-1}{4n}
$$
即 $2mn=-1$ 而这与假设条件 $m,n \ge 4$ 矛盾。
\end{solutions}